\documentclass{article}

\usepackage[dblblindworkshop, final]{neurips_2026}

\usepackage[utf8]{inputenc}
\usepackage[T1]{fontenc}
\usepackage{hyperref}
\usepackage{url}
\usepackage{booktabs}
\usepackage{amsfonts}
\usepackage{amsmath}
\usepackage{microtype}
\usepackage{xcolor}
\usepackage{graphicx}

\workshoptitle{30562 --- Machine Learning and Artificial Intelligence}
\title{Project Report Title}

\author{%
  Author One \\
  Bocconi University \\
  \texttt{author.one@studbocconi.it} \\
  \And
  Author Two \\
  Bocconi University \\
  \texttt{author.two@studbocconi.it} \\
  \And
  Author Three \\
  Bocconi University \\
  \texttt{author.three@studbocconi.it} \\
}



\begin{document}

\maketitle

\begin{abstract}
A concise summary of the problem, the approach, and the main findings.
The abstract should be self-contained and not exceed one paragraph.
\end{abstract}

\section{Introduction}

Motivate the problem you are solving.
What is the research question?
Why is it interesting or important?
Briefly summarise what you do and what you find.

\section{Background and Related Work}

Describe any prior work your project builds on.
Cite relevant papers and explain how your work relates to them.

\section{Method}

Describe your approach clearly and precisely.
Include any mathematical formulation, model architecture, or algorithm
that is central to your work.

\section{Experiments}

Describe your experimental setup: dataset(s), baselines, evaluation metrics,
and implementation details (model size, optimiser, hyperparameters).

\subsection{Results}

Present your main results using tables or figures.
Compare against baselines where applicable.

\subsection{Ablations}

If applicable, include ablation studies that isolate the contribution
of individual design choices.

\section{Conclusion}

Summarise what you did, what you found, and what the limitations are.
Optionally, suggest directions for future work.

\section*{References}

\small

% Use any consistent citation style.
% Example:
%
% [1] Author, A. \& Author, B. (Year). Title. \textit{Venue}.

\appendix

\section{Additional Results and Implementation Details}

Put here any supplementary figures, tables, proofs, or extended
experimental details that did not fit in the main paper.
The appendix has no page limit.

\end{document}
